\begin{statementbox}

	\textbf{\textcolor{CStatement}{条件}}
	\begin{description}[style=nextline, leftmargin=2.8em, labelsep=0.5em, itemsep=0.35em]
		\item[\textbf{\textcolor{CStatement}{条件一（原函数与图像）：}}] 设函数 $y=f(x)$ 的定义域为 $D$，其图像为曲线 $C$。
	\end{description}

	\textbf{\textcolor{CStatement}{结论}}
	\begin{description}[style=nextline, leftmargin=2.8em, labelsep=0.5em, itemsep=0.45em]
		\item[\textbf{\textcolor{CStatement}{结论一（横轴对称）：}}]
			\textbf{\textcolor{CStatement}{直观描述：}}
			图像关于 $x$ 轴对称后，纵坐标取相反数，横坐标不变。
			\[
				y=-f(x).
			\]

		\item[\textbf{\textcolor{CStatement}{结论二（纵轴对称）：}}]
			\textbf{\textcolor{CStatement}{直观描述：}}
			图像关于 $y$ 轴对称后，横坐标取相反数，纵坐标不变。
			\[
				y=f(-x).
			\]

		\item[\textbf{\textcolor{CStatement}{结论三（中心对称）：}}]
			\textbf{\textcolor{CStatement}{直观描述：}}
			图像关于原点对称后，横纵坐标均取相反数。
			\[
				y=-f(-x).
			\]

		\item[\textbf{\textcolor{CStatement}{推广结论（反函数对称）：}}]
			\textbf{\textcolor{CStatement}{直观描述：}}
			若 $f$ 存在反函数，则图像关于 $y=x$ 对称后对应反函数图像。
			\[
				y=f^{-1}(x).
			\]
	\end{description}
\end{statementbox}
